% Generated draft: requires mathematical and source review.
\documentclass[11pt]{article}
\usepackage[margin=25mm]{geometry}
\usepackage{fontspec}
\setmainfont{Latin Modern Roman}
\usepackage{amsmath,amssymb,mathtools,unicode-math}
\setmathfont{Latin Modern Math}
\usepackage{enumitem,needspace,xurl,hyperref}
\setlength{\parindent}{0pt}
\setlength{\parskip}{5pt}
\setlength{\emergencystretch}{3em}
\newcommand{\N}{\mathbb{N}}
\newcommand{\Z}{\mathbb{Z}}
\newcommand{\Q}{\mathbb{Q}}
\newcommand{\R}{\mathbb{R}}
\newcommand{\C}{\mathbb{C}}
\newcommand{\F}{\mathbb{F}}
\newcommand{\A}{\mathbb{A}}
\newcommand{\PP}{\mathbb{P}}
\newcommand{\E}{\mathbb{E}}
\newcommand{\eps}{\varepsilon}
\newcommand{\dd}{\,\mathrm{d}}
\begin{document}
\section*{UnsolvedMath ID: 800870}
\textbf{Source problem code:} catalogue entry
\subsection{Definitions and mathematical statement}
Let \(n\geq 1\) and let \(\Omega\subset\C^n\) be a domain, meaning a nonempty connected open set. A function \(u:\Omega\to[-\infty,\infty)\) is plurisubharmonic if it is upper semicontinuous and its restriction to each connected component of the intersection with a complex affine line is subharmonic or identically \(-\infty\). The source does not specify whether the globally identically \(-\infty\) function is admitted.

Let \(\varphi\) be an arbitrary plurisubharmonic function on \(\Omega\), and let \((\varphi_j)_{j\geq1}\) be bounded plurisubharmonic functions on \(\Omega\) satisfying
\[
\varphi_{j+1}(z)\leq\varphi_j(z),\qquad \lim_{j\to\infty}\varphi_j(z)=\varphi(z)\quad(z\in\Omega).
\]
Here bounded means bounded on \(\Omega\) for each individual \(j\); no bound uniform in \(j\) is imposed. Write
\[
\mu_j=(dd^c\varphi_j)^n
\]
for the Bedford--Taylor complex Monge--Amp\`ere measure of the bounded potential \(\varphi_j\). Fix one normalization of \(d^c\) throughout; the source does not specify that normalization.

For an explicit working interpretation of cluster points, use local weak convergence of positive Radon measures: \(\mu_{j_k}\to\mu\) if
\[
\lim_{k\to\infty}\int_\Omega f\,\dd\mu_{j_k}=\int_\Omega f\,\dd\mu
\quad\text{for every }f\in C_c(\Omega;\R).
\]
Thus, for a fixed approximating sequence, define
\[
\mathcal C((\varphi_j))=\left\{\mu\text{ a positive Radon measure on }\Omega:\ \mu_{j_k}\to\mu\text{ locally weakly for some }j_k\uparrow\infty\right\}.
\]
\textbf{Question (Bedford, as reported by Dinew, Guedj, and Zeriahi).} What can be said about the set of cluster points of \((dd^c\varphi_j)^n\) under these hypotheses?

The local weak topology above is an explicit working convention, not a topology specified in the supplied statement. The source also leaves unresolved whether the intended object is \(\mathcal C((\varphi_j))\) for a given approximating sequence or the union of these sets over all bounded decreasing approximations of a fixed \(\varphi\). No existence of cluster points or independence from the approximation is asserted.
\subsection{Short English statement}
For an arbitrary plurisubharmonic function, investigate the possible limiting complex Monge--Amp\`ere measures of bounded plurisubharmonic approximations decreasing pointwise to it. The source leaves the convergence topology and the role of the choice of approximation unspecified.
\subsection{Sources}
\begin{itemize}
\item Local UnsolvedMath registry, permanent ID \(800870\), source code \texttt{catalogue entry}, title \emph{Open Problems in Pluripotential Theory --- Question 1}. The supplied record labels the problem open; this label is not an independently verified assessment.
\item Dinew, Guedj, and Zeriahi, \emph{Open Problems in Pluripotential Theory} (2015), Question 1, attributed there to Bedford. Supplied source URL: \url{https://arxiv.org/abs/1511.00705}.
\item The supplied provenance records access on \texttt{2026-10-16} and extraction by \texttt{arxiv\_latex}. These are supplied provenance details, not a claim of independent source access or verification. Supplied upstream revision: \texttt{d371ee917ac9647fbb06cd6fc529deec6f3dda82}.
\end{itemize}
\end{document}
